\title[Israel 4: David--división]{\bfseries{Israel, el pueblo de Dios}}
\subtitle{Estudio 4. De David a la división del reino}
\author{}\institute{}\date{}
\titlegraphic{\includegraphics[width=0.85\linewidth,height=0.38\textheight,keepaspectratio]{portada.jpg}\\[-1pt]
  {\fontsize{5}{6}\selectfont\color{sepia}\textit{Nota.} Elaboración propia: 2~S 7:16 en texto hebreo consonántico (BHS).}}
\begin{document}
\frame{\titlepage}

% ===================================================
\section{Introducción}
\begin{frame}{¿Una promesa «para siempre»?}
\begin{columns}[T,onlytextwidth]
\begin{column}{0.52\textwidth}
\begin{cita}Y será firme tu casa y tu reino para siempre delante de ti, y tu trono será estable para siempre.\fuente{(2~S 7:16)}\end{cita}
\begin{cita}¿Qué parte tenemos nosotros con David? [\ldots] ¡Israel, a tus tiendas!\fuente{(1~R 12:16)}\end{cita}
\vspace{2pt}
\begin{clave}\footnotesize\destacar{El pacto} lleva la vocación de Éx 19:6 a un rey, una ciudad y una casa de oración «para todos los pueblos» (1~R 8:41--43).\end{clave}
\end{column}
\begin{column}{0.44\textwidth}\centering
\begin{tikzpicture}[node distance=7mm]
\node[bloqueazul,text width=4.6cm] (a) {Promesa\\[1pt]{\normalfont casa y trono «para siempre»}\\[1pt]{\normalfont\tiny 2~S 7:16}};
\node[bloque,text width=4.6cm,below=of a] (b) {\textbf{Cisma}\\[1pt]«¡Israel, a tus tiendas!»\\[1pt]{\tiny 1~R 12:16}};
\node[bloqueoro,text width=4.6cm,below=of b] (c) {\textbf{Juicio con límite}\\[1pt]afligida, «pero no para siempre»\\[1pt]{\tiny 1~R 11:39}};
\draw[flecha] (a) -- node[right=2pt,etiq]{menos de\\un siglo} (b);
\draw[flecha] (b) -- (c);
\end{tikzpicture}
\elabpropia
\end{column}
\end{columns}
\end{frame}

\begin{frame}{Los ejes de la serie en esta etapa}
\centering
\begin{tikzpicture}
 \foreach \x/\n/\t/\d/\r in {0/1/{Nación santa}/{¿Qué lugar tienen el rey, Sion y el templo en la vocación de Israel hacia las naciones?}/{2~S 6:14--18; 1~R 8:41--43; Sal 72:17},
   4.2/2/{Restauración}/{¿Qué límite pone Dios al juicio sobre la casa de David y el reino dividido?}/{2~S 7:14--16; 1~R 8:46--53; 11:39},
   8.4/3/{Proyección}/{¿Cómo esperan los profetas y Pablo al hijo de David y la reunificación?}/{Is 11; Jer 23:5--6; Ez 37; Ro 1:3--4; 15:12}}{
  \node[caja,anchor=north] (b\n) at (\x,0) {\parbox[t][3.4cm][t]{3.55cm}{\RaggedRight
    {\small\textbf{\textcolor{ladrillo}{\t}}}\\[5pt]\small \d\par\vfill{\scriptsize\color{sepia}(\r)}}};
  \node[circle,fill=cafe,text=piedra,font=\scriptsize\bfseries,inner sep=1.6pt] at (b\n.north west) {\n};}
 \draw[flecha] (b1.east) -- (b2.west);
 \draw[flecha] (b2.east) -- (b3.west);
\end{tikzpicture}
\elabpropia
\vspace{6pt}
\begin{notalam}[Sobre las citas]
Las citas bíblicas son traducción propia desde el texto hebreo (BHS) y griego (NA28). Numeración española [hebrea] cuando difiere.
\end{notalam}
\end{frame}

\begin{frame}{Tres sentidos de «Israel»}
\centering
\begin{tikzpicture}
\node[bloque,text width=4.4cm,minimum height=2.0cm,anchor=north west] (n) at (0.35,-0.95) {\textbf{\textcolor{ladrillo}{Israel (norte)}}\\[2pt]Diez tribus, centro en Efraín; capitales Siquem, Penuel, Tirsa y, desde Omri, Samaria\\[2pt]{\tiny\color{sepia}2~S 2:9; 1~R 12:16--20; Os 5:3}};
\node[bloqueoro,text width=4.4cm,minimum height=2.0cm,anchor=north east] (s) at (12.05,-0.95) {\textbf{\textcolor{ladrillo}{Judá (sur)}}\\[2pt]Judá y Benjamín; Jerusalén; casa de David\\[2pt]{\tiny\color{sepia}2~S 2:4; 1~R 12:21}};
\draw[draw=ladrillo,line width=0.9pt,rounded corners=5pt] (0,0) rectangle (12.4,-3.35);
\node[anchor=north,font=\small] at (6.2,-0.12) {\textbf{\textcolor{ladrillo}{Israel (pueblo)}}: las doce tribus, descendencia de Jacob {\scriptsize\color{sepia}(2~S 7:10; 1~R 8:62; Ro 11:1)}};
\draw[<->,draw=oro!90!ladrillo,line width=1pt,dashed] (n.east) -- node[above,etiq]{un pueblo,}node[below,etiq]{dos reinos} (s.west);
\end{tikzpicture}
\elabpropia
\vspace{4pt}
\begin{notalam}[«Todo Israel» cambia de alcance]
Todas las tribus en 1~R 8:62; solo las del norte en 1~R 12:16--20. Incluso bajo un solo rey se habla de «Israel y Judá» (2~S 5:5; 1~R 4:20). La pregunta vuelve en Ro 11:26, «todo Israel será salvo» (Estudio~8).
\end{notalam}
\end{frame}

\begin{frame}{Secuencia de la etapa}
\centering\vspace{8pt}
\begin{tikzpicture}[x=0.755cm]
\draw[line width=1.6pt,cafe,-{Latex[length=3mm]}] (0,0) -- (15.7,0);
\foreach \x/\t/\b in {0.6/{David rey en Hebrón}/{2~S 2:4; c. 1010},
                      2.7/{Rey de todo Israel; Jerusalén}/{2~S 5; c. 1003},
                      4.8/{Arca a Sion; pacto}/{2~S 6--7},
                      6.9/{Betsabé, Absalón}/{2~S 11--20},
                      9.0/{Salomón rey}/{1~R 1--2; c. 970},
                      11.1/{El templo}/{1~R 6--8; c. 966--959},
                      13.2/{División}/{1~R 12; c. 931/930},
                      15.1/{Sisac}/{1~R 14:25; c. 926/925}}{
  \fill[oro] (\x,0) circle (4pt);
  \node[above,font=\scriptsize\bfseries,text=ladrillo,text width=1.6cm,align=center] at (\x,0.2) {\t};
  \node[below,font=\tiny,text=sepia,text width=1.6cm,align=center] at (\x,-0.2) {\b};}
\draw[decorate,decoration={brace,amplitude=5pt,mirror},sepia] (0.4,-0.9) -- (7.1,-0.9) node[midway,below=5pt,font=\scriptsize\itshape]{David: un reino unido por pacto};
\draw[decorate,decoration={brace,amplitude=5pt,mirror},sepia] (8.8,-0.9) -- (13.4,-0.9) node[midway,below=5pt,font=\scriptsize\itshape]{Salomón: esplendor y quiebre};
\end{tikzpicture}
\notafig{Elaboración propia. Fechas aproximadas a.\,e.\,c. según Thiele (1983); otras cronologías bajan la división a c. 922.}
\vspace{10pt}
\begin{clave}\footnotesize Entre la promesa a David (2~S 7) y el cisma de Siquem (1~R 12) media menos de un siglo.\end{clave}
\end{frame}

\begin{frame}{La fecha de la división}
\centering
\begin{tikzpicture}[x=0.72cm]
\draw[line width=1.4pt,cafe,-{Latex[length=3mm]}] (0,0) -- (17.4,0);
\foreach \a in {935,930,925,920}{\draw[sepia] ({935-\a},0.1) -- ({935-\a},-0.1) node[below,font=\tiny,text=sepia]{\a};}
\node[hito] at (4.5,0) {};
\draw[oro!90!ladrillo,line width=0.8pt] (4.5,0.15) -- (4.5,0.75);
\node[bloqueazul,anchor=south,text width=2.9cm] at (4.0,0.75) {División, 931/930\\{\normalfont\tiny Thiele (1983); la que sigue esta serie}};
\node[hito] at (9.5,0) {};
\draw[oro!90!ladrillo,line width=0.8pt] (9.5,0.15) -- (9.5,0.75);
\node[bloqueoro,anchor=south,text width=2.9cm] at (10.0,0.75) {\textbf{Sisac en Judá e Israel}\\{\tiny 5.º año de Roboam (1~R 14:25), c. 926/925}};
\node[hito] at (13,0) {};
\draw[oro!90!ladrillo,line width=0.8pt] (13,-0.45) -- (13,-0.95);
\node[bloque,anchor=north,text width=3.3cm] at (13,-0.95) {\textbf{División, c. 922}\\{\tiny escuela de Albright}};
\end{tikzpicture}
\elabpropia
\vspace{4pt}
\begin{notalam}[Primer sincronismo externo: la lista de Karnak]
Sheshonq~I grabó en Karnak unos 150 topónimos del Levante (Gabaón, Meguido, Négueb). Jerusalén no figura en lo legible, lo que concuerda con el rescate de 1~R 14:25--26; la lista incluye ciudades del reino del norte, cuyo rey Jeroboam había sido asilado por el propio Sisac (1~R 11:40; Kitchen, 1986).
\end{notalam}
\end{frame}

\begin{frame}{¿800.000 en el norte, 400.000 en Judá?}
\begin{columns}[T,onlytextwidth]
\begin{column}{0.55\textwidth}
\footnotesize\renewcommand{\arraystretch}{1.25}
\begin{tabular}{>{\RaggedRight\arraybackslash}p{1.7cm}>{\RaggedRight\arraybackslash}p{4.5cm}}
\toprule
\color{ladrillo}\bfseries Fuente & \color{ladrillo}\bfseries Dato\\\midrule
2~S 24:9 & Guerreros, no habitantes: 800.000 Israel; 500.000 Judá\\
1~Cr 21:5 & 1.100.000 Israel; 470.000 Judá\\
\emph{Élef} & ¿«mil» o unidad militar? (Mendenhall, 1958)\\
Arqueología & Cientos de miles en todo el territorio en el s. \textsc{viii} (Broshi y Finkelstein, 1992)\\
\bottomrule
\end{tabular}
\elabpropia
\end{column}
\begin{column}{0.41\textwidth}
\begin{notalam}[\emph{Élef}: ¿mil o una unidad?]
El hebreo \heb{אֶלֶף} puede leerse como «mil» o como una unidad militar o clánica (Mendenhall, 1958). Según cómo se lea, las cifras de los censos cambian de orden de magnitud.
\end{notalam}
\vspace{4pt}
\begin{clave}\footnotesize Cifras estimadas; lo firme es la proporción: \destacar{un norte mayor y más rico} que Judá.\end{clave}
\end{column}
\end{columns}
\end{frame}

% ===================================================
\section{Arqueología}
\begin{frame}{Los dos reinos}
\begin{columns}[T,onlytextwidth]
\begin{column}{0.42\textwidth}\centering\includegraphics[width=\linewidth,height=0.74\textheight,keepaspectratio]{mapa_division.pdf}
\notafig{Elaboración propia (GSHHS/WDB, Basemap); ubicaciones aproximadas.}\end{column}
\begin{column}{0.54\textwidth}\scriptsize\RaggedRight
{\scriptsize\RaggedRight
\makebox[0.6cm][c]{\textcolor{ladrillo}{\rule{0.55ex}{0.55ex}}}~Capital\par
\makebox[0.6cm][c]{\textcolor{cafeosc}{\textbullet}}~Ciudad\par}
\par\smallskip
\textcolor{cafeosc}{\textbf{- - -}} Frontera aproximada: cambió con las guerras (1~R 15:16--22).\par\medskip
\begin{tcolorbox}[enhanced,colback=piedra,colframe=ladrillo,boxrule=0pt,leftrule=2.5pt,arc=0pt,left=5pt,right=5pt,top=3pt,bottom=3pt,fontupper=\scriptsize,before skip=2pt,after skip=4pt]
\textbf{\textcolor{ladrillo}{Israel (norte)}}: más extenso, fértil y poblado; capitales Siquem, Penuel, Tirsá y luego Samaria. Santuarios de Jeroboam en Betel y Dan (1~R 12:29).
\end{tcolorbox}
\begin{tcolorbox}[enhanced,colback=oro!12!white,colframe=oro!90!ladrillo,boxrule=0pt,leftrule=2.5pt,arc=0pt,left=5pt,right=5pt,top=3pt,bottom=3pt,fontupper=\scriptsize,before skip=2pt,after skip=4pt]
\textbf{\textcolor{ladrillo}{Judá (sur)}}: montañoso; Jerusalén, en la frontera con Benjamín, y la casa de David.
\end{tcolorbox}
Al este, Moab, donde la estela de Mesa (Dibón) nombra a Omri, rey de Israel.
\end{column}
\end{columns}
\end{frame}

\begin{frame}{La «casa de David» fuera de la Biblia}
\begin{columns}[T,onlytextwidth]
\begin{column}{0.48\textwidth}\footnotesize
\textbf{\textcolor{ladrillo}{Estela de Tel Dan}} (c. 840 a.\,e.\,c.)\par\smallskip
Aramea; Biran, 1993--1994; probablemente Hazael. Primera mención segura de David fuera de la Biblia (Biran y Naveh, 1993, 1995).\par\smallskip
\centering
\begin{tikzpicture}
\node[bloqueazul,text width=3.6cm] (h) at (0,0) {Rey arameo (¿Hazael?)};
\node[bloque,text width=2.4cm,minimum height=1.35cm,anchor=north] (i) at (-1.6,-0.95) {[Jo]ram, «rey de Israel»\\\heb{מלך ישראל}\\{\tiny = reino del norte}};
\node[bloqueoro,text width=2.4cm,minimum height=1.35cm,anchor=north] (d) at (1.6,-0.95) {rey de la «casa de David»\\\heb{ביתדוד}\\{\tiny = Judá}};
\draw[flecha] (h.south) -- ++(0,-0.2) -| (i.north);
\draw[flecha] (h.south) -- ++(0,-0.2) -| (d.north);
\end{tikzpicture}
\notafig{Elaboración propia a partir de Biran y Naveh (1993, 1995).}
\end{column}
\begin{column}{0.48\textwidth}\footnotesize
\textbf{\textcolor{ladrillo}{Estela de Mesa}} (c. 840 a.\,e.\,c.)
\begin{itemize}
\item Moabita; Dibón, 1868. Nombra a «Omri, rey de Israel» y a YHWH.
\item Línea 31: \emph{b[t]dwd} según Lemaire (1994).
\item Finkelstein, Na'aman y Römer (2019): quizá «Balac».
\end{itemize}
\begin{notalam}[Lectura debatida]
En Mesa la «casa de David» no es prueba firme. En Tel Dan, en cambio, un rey extranjero ve dos reinos: «Israel» y la «casa de David».
\end{notalam}
\end{column}
\end{columns}
\end{frame}

\begin{frame}{El siglo X a.\,e.\,c. en debate}
\centering\vspace{4pt}
\begin{tikzpicture}
 \foreach \x/\n/\t/\d/\r in {0/1/{Khirbet Qeiyafa}/{Ciudad fortificada con dos puertas sobre el valle de Elá, c. 1020--980. ¿Judaíta o de otra entidad?}/{Garfinkel et~al., 2016; Finkelstein y Fantalkin, 2012},
   3.15/2/{Ciudad de David}/{Estructura escalonada (¿Miló?) y Gran Estructura de Piedra (¿palacio de David?, 2~S 5:11).}/{E. Mazar, 2006; Finkelstein et~al., 2007},
   6.3/3/{Hazor, Meguido, Gezer}/{Puertas de seis cámaras (1~R 9:15): ¿Salomón (s. \textsc{x}) u omridas (s. \textsc{ix})?}/{Yadin, 1958; Finkelstein, 1996; A. Mazar, 2010},
   9.45/4/{Timna}/{Cobre de los s. \textsc{xi}--\textsc{x} sin edificios: el «sesgo arquitectónico».}/{Ben-Yosef, 2019}}{
  \node[caja,anchor=north] (b\n) at (\x,0) {\parbox[t][4.3cm][t]{2.65cm}{\RaggedRight
    {\small\textbf{\textcolor{ladrillo}{\t}}}\\[5pt]\footnotesize \d\par\vfill{\tiny\color{sepia}(\r)}}};
  \node[circle,fill=cafe,text=piedra,font=\scriptsize\bfseries,inner sep=1.6pt] at (b\n.north west) {\n};}
\end{tikzpicture}
\notafig{Elaboración propia a partir de los autores citados.}
\end{frame}

\begin{frame}{Dos lecturas del siglo X}
\centering\vspace{2pt}
\begin{tikzpicture}
\draw[{Latex[length=3mm]}-{Latex[length=3mm]},line width=1.4pt,cafe] (0,0) -- (12.2,0);
\node[bloque,anchor=north west,text width=4.3cm] at (0,-0.3) {\textbf{Cronología baja}\\[1pt]Finkelstein: David y Salomón, jefes de un reino pequeño};
\node[bloqueazul,anchor=north east,text width=4.3cm] at (12.2,-0.3) {Convencional modificada\\[1pt]{\normalfont A. Mazar y Garfinkel: un Judá organizado ya hacia el 1000}};
\node[bloqueoro,anchor=south,text width=4.4cm] at (6.1,0.3) {\textbf{Qeiyafa}: alguien fortificaba a gran escala en tiempos de David; lo que se discute es quién};
\node[hito] at (6.1,0) {};
\end{tikzpicture}
\notafig{Elaboración propia a partir de Finkelstein y Fantalkin (2012), A. Mazar (2010) y Garfinkel et~al. (2016).}
\vspace{6pt}
\begin{notalam}[El templo de Salomón y la arqueología]
Del templo de Salomón no hay restos excavados. Los paralelos más cercanos de su planta tripartita son templos sirios como el de Ain Dara (Monson, 2000). En las puertas de seis cámaras, el radiocarbono no ha cerrado la cuestión.
\end{notalam}
\end{frame}

% ===================================================
\section{David}
\begin{frame}{Hebrón: un pacto con el norte}
\begin{columns}[T,onlytextwidth]
\begin{column}{0.50\textwidth}
\centering{\large\heb{הִנְנוּ עַצְמְךָ וּבְשָׂרְךָ אֲנָחְנוּ}}\par\smallskip
\begin{cita}Henos aquí, hueso tuyo y carne tuya somos. [\ldots] Y el rey David hizo pacto con ellos en Hebrón delante de YHWH.\fuente{(2~S 5:1, 3)}\end{cita}
\footnotesize Una unidad sellada \destacar{por pacto} (\heb{בְּרִית}), que el norte querrá renegociar en Siquem (1~R 12:4).
\end{column}
\begin{column}{0.46\textwidth}\centering
\begin{tikzpicture}[node distance=8mm and 3mm]
\node[bloqueoro,text width=2.35cm,minimum height=1.35cm] (j) {\textbf{David}\\rey de «la casa de Judá»\\{\tiny 2~S 2:4}};
\node[bloque,text width=2.35cm,minimum height=1.35cm,right=of j] (n) {\textbf{Is-boset}\\rey de «todo Israel» (= norte)\\{\tiny 2~S 2:9}};
\node[bloqueazul,text width=5.2cm,below=10mm of $(j.south)!0.5!(n.south)$] (p) {Pacto en Hebrón\\{\normalfont un rey, dos entidades que se saben un pueblo}\\{\normalfont\tiny tras siete años y medio; 2~S 5:1--5}};
\draw[flecha] (j.south) -- (j.south |- p.north);
\draw[flecha] (n.south) -- (n.south |- p.north);
\end{tikzpicture}
\elabpropia
\end{column}
\end{columns}
\end{frame}

\begin{frame}{Un grito que vuelve}
\begin{cita}No tenemos parte en David ni heredad en el hijo de Isaí. ¡Cada uno a sus tiendas, Israel!\fuente{(2~S 20:1; cf. 1~R 12:16)}\end{cita}
\centering\vspace{4pt}
\begin{tikzpicture}[node distance=4mm]
\node[bloqueazul,text width=1.95cm,minimum height=1.45cm] (a) {Pacto\\{\normalfont\tiny Hebrón, 2~S 5:3}};
\node[bloque,text width=1.95cm,minimum height=1.45cm,right=of a] (b) {Absalón rey en Hebrón\\{\tiny 2~S 15:10}};
\node[bloque,text width=1.95cm,minimum height=1.45cm,right=of b] (c) {«Diez partes tenemos en el rey»\\{\tiny 2~S 19:43 [19:44]}};
\node[bloqueoro,text width=1.95cm,minimum height=1.45cm,right=of c] (d) {Seba toca el shofar\\{\tiny 2~S 20:1}};
\node[bloqueazul,text width=1.95cm,minimum height=1.45cm,right=of d] (e) {Siquem\\{\normalfont\tiny 1~R 12:16}};
\foreach \u/\v in {a/b,b/c,c/d,d/e}{\draw[flecha] (\u) -- (\v);}
\draw[decorate,decoration={brace,amplitude=5pt,mirror},sepia] (b.south west) ++(0,-0.15) -- ($(d.south east)+(0,-0.15)$) node[midway,below=6pt,font=\scriptsize\itshape]{tensión norte--sur bajo David};
\end{tikzpicture}
\elabpropia
\vspace{6pt}
\begin{clave}\footnotesize La unidad bajo David descansaba en un acuerdo entre dos entidades que se sabían un solo pueblo; el grito de Seba reaparece, con una línea más, en Siquem.\end{clave}
\end{frame}

\begin{frame}{Jerusalén y el arca en Sion}
\begin{columns}[T,onlytextwidth]
\begin{column}{0.48\textwidth}\centering
\begin{tikzpicture}
\node[bloqueoro,text width=1.4cm,minimum height=1.25cm] (j) at (-2.1,0) {\textbf{Judá}\\{\tiny Jos 15:8}};
\node[bloqueazul,text width=1.7cm,minimum height=1.25cm] (s) at (0,0) {Jebús / Sion\\{\normalfont\tiny plaza neutral}};
\node[bloque,text width=1.4cm,minimum height=1.25cm] (b) at (2.1,0) {\textbf{Benjamín}\\{\tiny Jos 18:16}};
\node[bloque,text width=1.7cm,fill=white,dashed] (si) at (-2.0,-2.1) {Silo, santuario de Efraín\\{\tiny Sal 78:60}};
\node[bloqueazul,text width=1.7cm] (ar) at (1.9,-2.1) {El arca sube a Sion\\{\normalfont\tiny 2~S 6; Sal 132:13--14}};
\draw[flecha] (si.east) -- node[above,etiq]{queda atrás} (ar.west);
\draw[flecha] (ar.north) to[out=90,in=-60] (s.south east);
\end{tikzpicture}
\elabpropia
\vspace{2pt}
{\footnotesize David danza «ceñido con un efod de lino» (2~S 6:14), ofrece y bendice en el nombre de YHWH (6:17--18).\par}
\end{column}
\begin{column}{0.48\textwidth}
\begin{cita}Desechó la tienda de José, y no escogió la tribu de Efraín, sino que escogió la tribu de Judá, el monte Sion, al cual amó [\ldots] y eligió a David su siervo.\fuente{(Sal 78:67--70)}\end{cita}
\begin{notalam}[Unidad y recelo]
Sion no pertenecía a ninguno de los dos bloques y nació como principio de unidad; para Efraín, que había perdido Silo (Sal 78:60; Jer 7:12), fue también fuente de resentimiento.
\end{notalam}
\end{column}
\end{columns}
\end{frame}

\begin{frame}{«Tú eres ese hombre»}
\begin{columns}[T,onlytextwidth]
\begin{column}{0.46\textwidth}\centering
\begin{tikzpicture}[node distance=2.6mm,every node/.append style={text width=5.2cm,font=\scriptsize}]
\node[bloque] (a) {\textbf{Pecado}: Betsabé y Urías {\tiny(2~S 11)}};
\node[bloque,below=of a] (b) {\textbf{Natán}: \heb{אַתָּה הָאִישׁ} {\tiny(12:7)}};
\node[bloqueoro,below=of b] (c) {\textbf{Confesión}: \heb{חָטָאתִי לַיהוָה} {\tiny(12:13)}};
\node[bloque,below=of c] (d) {\textbf{Perdón y consecuencia}: «no se apartará jamás la espada de tu casa» {\tiny(12:10)}};
\node[bloqueazul,below=of d] (e) {Salomón, \heb{יְדִידְיָהּ}, «amado de YH» {\normalfont\tiny(12:25)}};
\node[bloque,below=of e] (f) {Absalón se proclama rey en Hebrón {\tiny(15:10)}};
\foreach \u/\v in {a/b,b/c,c/d,d/e,e/f}{\draw[flecha] (\u) -- (\v);}
\end{tikzpicture}
\elabpropia
\end{column}
\begin{column}{0.50\textwidth}
\centering{\large\heb{לֵב טָהוֹר בְּרָא־לִי אֱלֹהִים}}\par\smallskip
\begin{cita}Crea en mí, oh Dios, un corazón limpio, y renueva un espíritu firme dentro de mí.\fuente{(Sal 51:10 [51:12])}\end{cita}
\begin{notalam}[Hebreo]
\heb{בָּרָא}: verbo de Gn 1:1. \heb{נָכוֹן}, «firme», como el trono de 2~S 7:16 (observación propia): la estabilidad de la casa pasa por un corazón renovado.
\end{notalam}
\begin{clave}La era de Arauna (2~S 24) será el monte Moriah del templo (2~Cr 3:1; cf. Gn 22).\end{clave}
\end{column}
\end{columns}
\end{frame}

% ===================================================
\section{El pacto}
\begin{frame}{Una casa para David}
\fotolado{panel_2s7.jpg}{Elaboración propia: 2~S 7:13--14 en texto hebreo consonántico (BHS).}
\begin{columns}[T,onlytextwidth]
\begin{column}{0.57\textwidth}
\begin{cita}Él edificará casa a mi nombre, y yo afirmaré para siempre el trono de su reino. Yo le seré a él padre, y él me será a mí hijo [\ldots]; mi misericordia no se apartará de él.\fuente{(2~S 7:13--15)}\end{cita}
\centering
\begin{tikzpicture}[node distance=4mm]
\node[bloque,text width=1.55cm] (a) {\textbf{palacio}\\{\tiny 7:2, «casa de cedro»}};
\node[bloque,text width=1.55cm,right=of a] (b) {\textbf{templo}\\{\tiny 7:5, «¿tú me edificarás casa?»}};
\node[bloqueazul,text width=1.55cm,right=of b] (c) {dinastía\\{\normalfont\tiny 7:11, «YHWH te hará casa»}};
\draw[flecha] (a) -- (b); \draw[flecha] (b) -- (c);
\node[etiq,above=2pt of b.north] {\heb{בַּיִת} (\emph{báyit}): tres sentidos de «casa»};
\end{tikzpicture}
\elabpropia
\vspace{2pt}
\begin{clave}\footnotesize David quería edificarle casa a Dios; es Dios quien se la edifica a él: «firme tu casa [\ldots] para siempre» (7:16).\end{clave}
\end{column}
\begin{column}{0.41\textwidth}\end{column}
\end{columns}
\end{frame}

\begin{frame}{Cuatro claves de 2 S 7}
\centering
\begin{tikzpicture}
 \foreach \x/\n/\t/\d/\r in {0/1/{Firme}/{\heb{נֶאְמַן} (raíz \emph{'mn}) y \heb{נָכוֹן} (raíz \emph{kun}): la promesa descansa en la fidelidad de Dios.}/{2~S 7:16; Is 55:3},
   3.15/2/{Para siempre}/{\heb{עַד־עוֹלָם}: tiempo sin límite a la vista, reforzado por el sol y la luna.}/{Sal 89:36--37 [89:37--38]},
   6.3/3/{Hijo}/{Adopción real: «yo le seré padre». Israel es también hijo (Éx 4:22).}/{Sal 2:7; Heb 1:5; 2~Co 6:18},
   9.45/4/{Disciplina}/{Castigo con «vara de hombres», pero la \emph{jésed} no se aparta.}/{2~S 7:14--15; Sal 89:30--33}}{
  \node[caja,anchor=north] (b\n) at (\x,0) {\parbox[t][3.0cm][t]{2.65cm}{\RaggedRight
    {\small\textbf{\textcolor{ladrillo}{\t}}}\\[4pt]\footnotesize \d\par\vfill{\tiny\color{sepia}(\r)}}};
  \node[circle,fill=cafe,text=piedra,font=\scriptsize\bfseries,inner sep=1.6pt] at (b\n.north west) {\n};}
\end{tikzpicture}
\elabpropia
\vspace{4pt}
\begin{notalam}[Un pacto sin la palabra «pacto»]
\emph{Berit} no aparece en 2~S 7, pero otros textos lo llaman «pacto perpetuo» (2~S 23:5; Sal 89:3 [89:4]; 132:12; Jer 33:21). La forma condicional (1~R 2:4; Sal 132:12) no anula la promesa: afecta a cada rey, no a la casa.
\end{notalam}
\end{frame}

\begin{frame}{«Él me será hijo»: del hebreo a Pablo}
\centering
\begin{tikzpicture}[node distance=0.9cm]
\node[caja,text width=3.3cm,align=center] (h) {\textbf{Hebreo (2~S 7:14)}\\[3pt]\heb{וְהוּא יִהְיֶה־לִּי לְבֵן}\\[2pt]el hijo de David,\\hijo de Dios};
\node[caja,text width=3.3cm,align=center,right=of h] (g) {\textbf{LXX}\\[3pt]\grc{αὐτὸς ἔσται μοι εἰς υἱόν}\\[2pt]citado en Heb 1:5};
\node[caja,text width=3.3cm,align=center,right=of g] (p) {\textbf{Pablo}\\[3pt]\grc{εἰς υἱοὺς καὶ θυγατέρας}\\[2pt]2~Co 6:18; Ro 1:3--4};
\draw[-{Latex[length=2.5mm]},oro!90!ladrillo,line width=1.4pt] (h) -- (g);
\draw[-{Latex[length=2.5mm]},oro!90!ladrillo,line width=1.4pt] (g) -- (p);
\end{tikzpicture}
\elabpropia
\vspace{6pt}
\scriptsize\renewcommand{\arraystretch}{1.25}
\begin{tabular}{>{\RaggedRight\arraybackslash}p{2.3cm}>{\RaggedRight\arraybackslash}p{4.2cm}>{\RaggedRight\arraybackslash}p{4.2cm}}
\toprule
& \color{ladrillo}\bfseries 2~S 7 & \color{ladrillo}\bfseries 1~Cr 17\\\midrule
Descendiente & «que saldrá de tus entrañas» (7:12) & «uno de tus hijos» (17:11)\\
Castigo & «si hace mal, lo castigaré» (7:14) & omitido (17:13)\\
Reino & «\emph{tu} casa y \emph{tu} reino» (7:16) & «\emph{mi} casa y \emph{mi} reino» (17:14)\\
\bottomrule
\end{tabular}
\elabpropia
\end{frame}

% ===================================================
\section{Salomón}
\begin{frame}{Salomón: un corazón que escucha}
\begin{columns}[T,onlytextwidth]
\begin{column}{0.48\textwidth}\footnotesize
\begin{itemize}
\item \heb{לֵב שֹׁמֵעַ}: «un corazón que escucha» (1~R 3:9), con el verbo del Shemá.
\item «Judá e Israel eran muchos, como la arena del mar» (1~R 4:20; Gn 22:17).
\item «Cada uno debajo de su vid y de su higuera» (1~R 4:25 [5:5]; Mi 4:4).
\item Dedicación del templo en Sucot, «el mes séptimo» (8:2).
\end{itemize}
\begin{notalam}[Las columnas]
\heb{יָכִין} (\emph{Yajín}, «él establecerá», de la raíz \emph{kun} de 2~S 7:13, 16) y \heb{בֹּעַז} (\emph{Bóaz}), ante la entrada del templo (1~R 7:21; Cogan, 2001).
\end{notalam}
\end{column}
\begin{column}{0.48\textwidth}\centering
\begin{tikzpicture}[node distance=9mm]
\node[bloque,text width=4.6cm] (e) {\textbf{Éxodo de Egipto}\\{\tiny «la casa de YHWH» como meta del Éxodo (Éx 15:17)}};
\node[bloqueazul,text width=4.6cm,below=of e] (t) {Templo de Salomón\\{\normalfont\tiny cuarto año de Salomón, c. 966}};
\draw[flecha] (e) -- node[right=3pt,etiq]{480 años\\(1~R 6:1)} (t);
\node[bloqueoro,text width=4.6cm,below=7mm of t] (j) {\textbf{Josefo}: 592 años\\{\tiny \emph{Ant.} 8.3.1}};
\draw[dashed,draw=sepia,line width=0.7pt] (t) -- node[right=3pt,etiq]{otro cómputo} (j);
\end{tikzpicture}
\elabpropia
\end{column}
\end{columns}
\end{frame}

\begin{frame}{La gloria llena la casa}
\begin{columns}[T,onlytextwidth]
\begin{column}{0.50\textwidth}
\begin{cita}La nube llenó la casa de YHWH, y los sacerdotes no pudieron permanecer para ministrar [\ldots], porque la gloria de YHWH había llenado la casa.\fuente{(1~R 8:10--11)}\end{cita}
\begin{cita}Los cielos de los cielos no te pueden contener; ¡cuánto menos esta casa!\fuente{(1~R 8:27)}\end{cita}
\end{column}
\begin{column}{0.46\textwidth}\centering
\begin{tikzpicture}[node distance=8mm]
\node[bloque,text width=4.3cm] (a) {\textbf{Tabernáculo}\\la nube cubre la tienda\\{\tiny Éx 40:34--35}};
\node[bloqueazul,text width=4.3cm,below=of a] (b) {Templo\\{\normalfont la nube llena la casa}\\{\normalfont\tiny 1~R 8:10--11}};
\node[bloqueoro,text width=4.3cm,below=of b] (c) {\textbf{Nombre y oración}\\Dios pone allí su nombre; hacia allí se ora\\{\tiny 1~R 8:29--30}};
\draw[flecha] (a) -- node[right=3pt,etiq]{lo mismo sucede} (b);
\draw[flecha] (b) -- node[right=3pt,etiq]{no lo encierra} (c);
\end{tikzpicture}
\elabpropia
\end{column}
\end{columns}
\end{frame}

\begin{frame}{Una oración que prevé el exilio}
\centering{\large\heb{וְהֵשִׁיבוּ אֶל־לִבָּם בָּאָרֶץ אֲשֶׁר נִשְׁבּוּ־שָׁם וְשָׁבוּ}}\par\medskip
\begin{cita}Si pecaren contra ti, porque no hay hombre que no peque, [\ldots] y ellos volvieren en sí en la tierra adonde fueron llevados cautivos, y se convirtieren y oraren a ti [\ldots] con el rostro hacia su tierra, [\ldots] tú oirás en los cielos [\ldots] y les harás hallar misericordia.\fuente{(1~R 8:46--48, 50)}\end{cita}
\vspace{6pt}
\begin{clave}\footnotesize Siete súplicas (8:31--53); desde la segunda se pide «vuélvelos a la tierra» (8:34). \destacar{En la hora de su mayor gloria, Israel ya ora por su exilio.}\end{clave}
\end{frame}

\begin{frame}{Volver de corazón, volver a la tierra}
\begin{columns}[T,onlytextwidth]
\begin{column}{0.52\textwidth}\centering
\begin{tikzpicture}[every node/.append style={text width=2.45cm,minimum height=1.05cm}]
\node[bloque] (p) at (0,1.75) {\textbf{Pecado}\\{\tiny «no hay hombre que no peque»}};
\node[bloque] (c) at (3.2,1.75) {\textbf{Cautiverio}\\\heb{שָׁבָה} \emph{shavá}};
\node[bloqueoro] (v) at (3.2,-0.1) {\textbf{Vuelven en sí}\\\heb{שׁוּב} \emph{shuv}};
\node[bloque] (o) at (3.2,-1.95) {\textbf{Oran hacia}\\{\tiny la tierra, la ciudad, la casa}};
\node[bloqueazul] (m) at (0,-1.95) {Misericordia y regreso\\{\normalfont\tiny Dt 30:1--5}};
\draw[flecha] (p) -- (c); \draw[flecha] (c) -- (v); \draw[flecha] (v) -- (o); \draw[flecha] (o) -- (m);
\draw[flecha,dashed] (m) -- node[left,etiq,text width=1.3cm,minimum height=0pt]{1~R 8:34,\\46--50} (p);
\end{tikzpicture}
\elabpropia
\end{column}
\begin{column}{0.44\textwidth}
\begin{notalam}[Juego de palabras en hebreo]
\heb{שׁוּב} (\emph{shuv}, «volver»: \emph{ve-heshivu el libbam}, \emph{ve-shavu}) se entrelaza con \heb{שָׁבָה} (\emph{shavá}, «llevar cautivo»: \emph{nishbu}). El sonido dice lo que el texto enseña: el regreso del corazón precede al regreso a la tierra.
\end{notalam}
\begin{clave}La súplica se funda en el Éxodo: «ellos son tu pueblo y tu heredad, que tú sacaste de Egipto» (1~R 8:51; Dt 4:20).\end{clave}
\end{column}
\end{columns}
\end{frame}

\begin{frame}{Nación santa: vivida, quebrada, preservada}
\centering\vspace{2pt}
\begin{tikzpicture}[node distance=5mm]
\node[bloque,text width=3.45cm,minimum height=3.9cm,align=left,anchor=north] (a) at (0,0) {\textbf{\textcolor{ladrillo}{Vivida}}\par\smallskip
David con efod, ofrece y bendice (2~S 6:14--18). «Sacerdote [\ldots] según Melquisedec» (Sal 110:4). Salomón ora por el extranjero (1~R 8:41--43). La reina de Sabá bendice a YHWH (1~R 10:9). «Benditas serán en él todas las naciones» (Sal 72:17; Gn 12:3).};
\node[bloqueoro,text width=3.45cm,minimum height=3.9cm,align=left,anchor=north] (b) at (4.25,0) {\textbf{\textcolor{ladrillo}{Quebrada}}\par\smallskip
La leva (\emph{mas}) convierte al pueblo en un nuevo Egipto. Quien oró por las naciones erige altares a sus dioses (1~R 11:7). Jeroboam desliga el sacerdocio de la elección (1~R 12:31).};
\node[bloque,text width=3.45cm,minimum height=3.9cm,align=left,anchor=north] (c) at (8.5,0) {\textbf{\textcolor{ladrillo}{Preservada}}\par\smallskip
Profetas sobre el rey: Natán, Gad, Ahías, Semaías. Arrepentimiento (Sal 51; 1~R 8:47). Remanente en ambos reinos: levitas en Judá (2~Cr 11) y los siete mil del norte (1~R 19:18; Ro 11:4).};
\draw[flecha] (a) -- (b); \draw[flecha] (b) -- (c);
\end{tikzpicture}
\elabpropia
\vspace{4pt}
\begin{clave}\footnotesize «Mi casa será llamada casa de oración para todos los pueblos» (Is 56:7).\end{clave}
\end{frame}

% ===================================================
\section{La división}
\begin{frame}{Cuatro causas de la división}
\centering
\begin{tikzpicture}
 \foreach \x/\n/\t/\d/\r in {0/1/{Idolatría}/{Altares a Quemos y Moloc frente a Jerusalén; «sus mujeres desviaron su corazón».}/{1~R 11:1--13},
   3.15/2/{Trabajo forzado}/{\heb{מַס} (\emph{mas}) sobre «todo Israel»; Jeroboam, sobre «la casa de José».}/{1~R 5:13 [5:27]; 11:28},
   6.3/3/{Profecía}/{Ahías de Silo rasga el manto en doce pedazos.}/{1~R 11:29--39},
   9.45/4/{Necedad}/{Roboam en Siquem: «yo os castigaré con escorpiones».}/{1~R 12:1--15}}{
  \node[caja,anchor=north] (b\n) at (\x,0) {\parbox[t][2.6cm][t]{2.65cm}{\RaggedRight
    {\small\textbf{\textcolor{ladrillo}{\t}}}\\[4pt]\footnotesize \d\par\vfill{\tiny\color{sepia}(\r)}}};
  \node[circle,fill=cafe,text=piedra,font=\scriptsize\bfseries,inner sep=1.6pt] at (b\n.north west) {\n};}
 \node[bloqueazul,text width=5.2cm] (d) at ($(b2.south)!0.5!(b3.south)+(0,-1.0)$) {División del reino {\normalfont\tiny (1~R 12, c. 931/930)}};
 \foreach \n in {1,2,3,4}{\draw[flecha] (b\n.south) -- (d);}
\end{tikzpicture}
\elabpropia
\vspace{2pt}
\begin{clave}\footnotesize En el fondo, la ley del rey quebrantada: caballos de Egipto, mujeres, oro (Dt 17:16--17; 1~R 10:14, 26--29; 11:3). A ello se suma una tensión norte--sur que venía de antes (2~S 19--20).\end{clave}
\end{frame}

\begin{frame}{El rey se hace faraón}
\centering\vspace{2pt}
\begin{tikzpicture}[node distance=6mm]
\node[bloque,text width=4.6cm,minimum height=1.6cm] (e) {\textbf{Egipto}\\\heb{שָׂרֵי מִסִּים}, «capataces de la leva»\\{\tiny Éx 1:11}};
\node[bloqueoro,text width=4.6cm,minimum height=1.6cm,right=22mm of e] (s) {\textbf{Salomón}\\\heb{מַס} sobre «todo Israel»; Jeroboam, sobre «la casa de José»\\{\tiny 1~R 5:13 [5:27]; 11:28}};
\draw[<->,draw=oro!90!ladrillo,line width=1pt] (e) -- node[above,etiq]{la misma}node[below,etiq]{palabra} (s);
\node[bloqueazul,text width=6cm,below=9mm of $(e.south)!0.5!(s.south)$] (y) {«Tu padre agravó nuestro yugo»\\{\normalfont\tiny 1~R 12:4}};
\draw[flecha] (s.south) |- (y.east);
\end{tikzpicture}
\elabpropia
\vspace{4pt}
\begin{notalam}[¿Leva o servidumbre?]
1~R 9:20--22 reserva la servidumbre permanente a los cananeos; a los israelitas, una prestación por turnos (Cogan, 2001). Para el norte, esa distinción no bastó.
\end{notalam}
\end{frame}

\begin{frame}{Diez pedazos para Jeroboam}
\fotolado{panel_1r11.jpg}{Elaboración propia: 1~R 11:39 en texto hebreo consonántico (BHS).}
\begin{columns}[T,onlytextwidth]
\begin{column}{0.57\textwidth}
\begin{cita}[fontupper=\footnotesize\itshape]A ti te daré diez tribus; y él tendrá una tribu por amor a David mi siervo, y por amor a Jerusalén [\ldots], para que mi siervo David tenga lámpara todos los días delante de mí.\fuente{(1~R 11:31--32, 36)}\end{cita}
\begin{cita}[fontupper=\footnotesize\itshape]Y yo afligiré a la descendencia de David a causa de esto, \textbf{pero no para siempre}.\fuente{(1~R 11:39)}\end{cita}
\begin{notalam}[Hebreo y aritmética]
«Pero no todos los días» traduce \heb{אַךְ לֹא כָל־הַיָּמִים}; la «lámpara» (\heb{נִיר}) es la continuidad de la casa de David. 10 + 1 no son 12: ¿Benjamín, Simeón absorbido, Leví sin territorio? (1~R 12:21).
\end{notalam}
\end{column}
\begin{column}{0.41\textwidth}\end{column}
\end{columns}
\end{frame}

\begin{frame}{Siquem: «¡A tus tiendas, Israel!»}
\centering{\large\heb{מַה־לָּנוּ חֵלֶק בְּדָוִד}}\par\smallskip
\begin{cita}¿Qué parte tenemos nosotros con David? No tenemos heredad en el hijo de Isaí. ¡Israel, a tus tiendas! ¡Provee ahora en tu casa, David!\fuente{(1~R 12:16)}\end{cita}
\vspace{2pt}
\begin{tikzpicture}[node distance=5mm]
\node[bloque,text width=2.5cm,minimum height=1.35cm] (a) {\textbf{Siquem}\\lugar del pacto; el norte pone condiciones\\{\tiny Jos 24; 1~R 12:4}};
\node[bloqueoro,text width=2.5cm,minimum height=1.35cm,right=of a] (r) {\textbf{Roboam}\\«con escorpiones»; «la causa venía de YHWH»\\{\tiny 12:14--15}};
\node[bloqueazul,text width=3.2cm,anchor=south west] (n) at ($(r.east)+(0.7,0.08)$) {Israel (norte): Jeroboam rey de «todo Israel» {\normalfont\tiny(12:20)}};
\node[bloqueoro,text width=3.2cm,anchor=north west,font=\scriptsize\bfseries] (s) at ($(r.east)+(0.7,-0.08)$) {Judá y Benjamín siguen a David {\normalfont\tiny(12:21)}};
\draw[flecha] (a) -- (r);
\draw[flecha] (r.east) -- (n.west); \draw[flecha] (r.east) -- (s.west);
\end{tikzpicture}
\elabpropia
\vspace{4pt}
\begin{clave}\footnotesize Semaías detiene a Roboam: «no peleéis contra vuestros \textbf{hermanos}» (12:24). Consumada la ruptura, Israel (norte) y Judá siguen siendo hermanos.\end{clave}
\end{frame}

\begin{frame}{Los becerros de Betel y Dan}
\begin{columns}[T,onlytextwidth]
\begin{column}{0.48\textwidth}
\centering{\small\heb{הִנֵּה אֱלֹהֶיךָ יִשְׂרָאֵל}}\par\smallskip
\begin{cita}He aquí tus dioses, oh Israel, los cuales te hicieron subir de la tierra de Egipto.\fuente{(1~R 12:28; cf. Éx 32:4)}\end{cita}
\footnotesize
\begin{itemize}
\item El norte nace reeditando el pecado del Sinaí.
\item Sacerdotes no levitas; fiesta del mes octavo, copia de Sucot «inventada de su corazón» (12:31--33).
\end{itemize}
\end{column}
\begin{column}{0.48\textwidth}\centering
\begin{tikzpicture}[node distance=9mm]
\node[bloque,text width=4.6cm] (n) {\textbf{Israel (norte)}\\santuarios de Betel y Dan};
\node[bloqueoro,text width=4.6cm,below=of n] (s) {\textbf{Judá (sur)}\\culto de Jerusalén};
\draw[flecha] (n) -- node[right=3pt,etiq,text width=2.6cm,align=left]{sacerdotes, levitas y fieles «de todas las tribus»} (s);
\end{tikzpicture}
\notafig{Elaboración propia a partir de 2~Cr 11:13--17; 15:9; 30.}
\begin{notalam}[¿Ídolos o pedestales?]
Muchos especialistas ven en los becerros pedestales del Dios invisible, como los querubines (Cogan, 2001); el narrador los llama «el pecado de Jeroboam» (1~R 15:30).
\end{notalam}
\end{column}
\end{columns}
\end{frame}

% ===================================================
\section{Tradición judía}
\begin{frame}{El renuevo de David en la Amidá}
\centering{\large\heb{אֶת צֶמַח דָּוִד עַבְדְּךָ מְהֵרָה תַצְמִיחַ}}\par\smallskip
\begin{cita}Haz brotar pronto el renuevo de David tu siervo, y exalta su poder con tu salvación, porque tu salvación esperamos todo el día.\fuente{(Amidá, bendición 15)}\end{cita}
\begin{columns}[T,onlytextwidth]
\begin{column}{0.48\textwidth}\footnotesize
\textbf{\textcolor{cafe}{Origen bíblico}}: \emph{tsémaj} (Jer 23:5; 33:15); «haré brotar el poder de David» (Sal 132:17). Bendición 14: «el trono de David establece pronto» en Jerusalén.\par\medskip
\textbf{\textcolor{cafe}{Birkat ha-Mazón}}: «sobre el reino de la casa de David tu ungido» (Berajot 48b). En Sucot: «levante la cabaña caída de David» (Am 9:11).
\end{column}
\begin{column}{0.48\textwidth}
\begin{notalam}[Jerusalén y David, unidos]
En el rito de la tierra de Israel conservado en la Genizá de El Cairo, la bendición de David estaba unida a la de Jerusalén (Elbogen, 1993): la restauración de la ciudad y la del trono davídico se piden juntas.
\end{notalam}
\end{column}
\end{columns}
\end{frame}

\begin{frame}{David y Jeroboam en el Talmud}
\centering\vspace{2pt}
\begin{tikzpicture}
\node[bloque,text width=5.3cm,minimum height=4.6cm,align=left,anchor=north] (d) at (0,0) {\textbf{\textcolor{ladrillo}{David}}\par\smallskip
\textbf{El arpa}: el viento del norte lo despertaba a medianoche (Berajot 3b).\par\smallskip
\textbf{El pecado}: «para enseñar el arrepentimiento al individuo» (Avodá Zará 4b--5a); otra opinión niega que pecara (Shabat 56a).\par\smallskip
\textbf{La esperanza}: «David, rey de Israel, vive y permanece» (Rosh Hashaná 25a); «otro David» (Sanedrín 98b).};
\node[bloqueoro,text width=5.3cm,minimum height=4.6cm,align=left,anchor=north] (j) at (6.4,0) {\textbf{\textcolor{ladrillo}{Jeroboam}}\par\smallskip
\textbf{El juicio}: sin parte en el mundo venidero (Mishná, Sanedrín 10:2).\par\smallskip
\textbf{El diálogo}: Dios le ofrece pasear en el Edén si se arrepiente. «¿Quién irá delante?» «El hijo de Isaí». «Entonces no quiero» (Sanedrín 102a).\par\smallskip
\textbf{La lectura}: un orgullo incapaz de aceptar el lugar de la casa de David.};
\draw[<->,draw=oro!90!ladrillo,line width=1pt] (d.east) -- (j.west);
\end{tikzpicture}
\elabpropia
\end{frame}

\begin{frame}{Jerusalén en el rito}
\begin{columns}[T,onlytextwidth]
\begin{column}{0.46\textwidth}\centering
\begin{tikzpicture}
\fill[piedra,draw=ladrillo,line width=0.6pt] (0,0) circle (2.45);
\fill[oro!14!white,draw=oro!90!ladrillo,line width=0.6pt] (0,-0.6) circle (1.75);
\fill[celeste!35!white,draw=ladrillo,line width=0.6pt] (0,-1.15) circle (1.12);
\fill[ladrillo] (0,-1.65) circle (0.55);
\node[etiq] at (0,1.78) {fuera de la tierra:\\\textbf{hacia la tierra}};
\node[etiq] at (0,0.52) {en la tierra:\\\textbf{hacia Jerusalén}};
\node[etiq] at (0,-0.55) {en Jerusalén:\\\textbf{hacia el templo}};
\node[font=\tiny\bfseries,text=white,align=center] at (0,-1.65) {santí-\\simo};
\end{tikzpicture}
\notafig{Elaboración propia a partir de Berajot 30a; cf. 1~R 8:44, 48.}
\end{column}
\begin{column}{0.50\textwidth}\footnotesize
\textbf{\textcolor{cafe}{Yom Kipur}}: el orden de la confesión se discute desde 1~R 8:47 (Yomá 36b).\par\smallskip
\textbf{\textcolor{cafe}{Tisha be-Av}}: caída de los dos templos (Mishná, Taanit 4:6).
\begin{notalam}[Antes del Nuevo Testamento]
4Q174 aplica 2~S 7 al «Renuevo de David» (\heb{צמח דויד}); \emph{Sal. Salom.} 17:21 pide «el hijo de David»; el Targum de Is 11:1 lee «el Mesías». Josefo narra la etapa en \emph{Ant.} 7--8.
\end{notalam}
\end{column}
\end{columns}
\end{frame}

% ===================================================
\section{Restauración}
\begin{frame}{El límite del juicio}
\centering\footnotesize\renewcommand{\arraystretch}{1.25}
\begin{tabular}{>{\RaggedRight\arraybackslash}p{1.9cm}>{\RaggedRight\arraybackslash}p{4.7cm}>{\RaggedRight\arraybackslash}p{4.4cm}}
\toprule
\color{ladrillo}\bfseries Texto & \color{ladrillo}\bfseries Promesa & \color{ladrillo}\bfseries Proyección\\\midrule
2~S 7:14--16 & La \emph{jésed} no se aparta; trono «para siempre» & Is 55:3; Lc 1:32--33\\
Sal 132:11, 17 & Juramento; renuevo y lámpara & Jer 23:5; Hch 2:30\\
1~R 8:46--53 & Regreso del corazón y de la tierra & Dt 30:1--5; Dn 9:4--19\\
1~R 11:36, 39 & Lámpara; afligir «pero no para siempre» & Am 9:11; Os 3:5; Ez 37:22--25\\
\bottomrule
\end{tabular}
\elabpropia
\vspace{4pt}
\begin{cita}Señor, ¿dónde están tus antiguas misericordias, que juraste a David por tu fidelidad?\fuente{(Sal 89:49 [89:50])}\end{cita}
{\footnotesize\itshape\color{sepia}El Sal 89 no resuelve la tensión entre la promesa y la corona caída; la presenta ante Dios.}
\end{frame}

\begin{frame}{Isaías: el tronco de Isaí}
\begin{columns}[T,onlytextwidth]
\begin{column}{0.49\textwidth}
\centering{\heb{וְיָצָא חֹטֶר מִגֵּזַע יִשָׁי}}\par\smallskip
\begin{cita}Saldrá una vara del tronco de Isaí [\ldots]. La raíz de Isaí, puesta por pendón a los pueblos, será buscada por las naciones.\fuente{(Is 11:1, 10)}\end{cita}
\begin{notalam}[\heb{גֵּזַע}, «tocón»]
El árbol ha sido cortado, y el profeta nombra a Isaí como si la dinastía tuviera que empezar de nuevo. Pablo cita Is 11:10 en Ro 15:12 (LXX, \grc{ὁ ἀνιστάμενος}, «el que se levanta»).
\end{notalam}
\end{column}
\begin{column}{0.47\textwidth}
\begin{cita}[fontupper=\scriptsize\itshape]Juntará los desterrados de Israel, y reunirá los esparcidos de Judá [\ldots]. Efraín no tendrá envidia de Judá, ni Judá afligirá a Efraín.\fuente{(Is 11:12--13)}\end{cita}
\footnotesize
\begin{itemize}
\item Is 9:6--7 [9:5--6]: el niño en el trono de David; luz primero sobre Galilea, el norte deportado.
\item Is 55:3: «las misericordias firmes de David» (Hch 13:34).
\end{itemize}
\end{column}
\end{columns}
\end{frame}

\begin{frame}{Jeremías: el renuevo justo}
\begin{columns}[T,onlytextwidth]
\begin{column}{0.50\textwidth}
\centering{\heb{וַהֲקִמֹתִי לְדָוִד צֶמַח צַדִּיק}}\par\smallskip
\begin{cita}Levantaré a David renuevo justo [\ldots]. En sus días será salvo \textbf{Judá}, e \textbf{Israel} habitará confiado; y este será su nombre: YHWH, justicia nuestra.\fuente{(Jer 23:5--6)}\end{cita}
\begin{notalam}[Texto]
Jer 33:14--26 falta en la LXX, más breve en Jeremías; se sigue el TM. Su contenido coincide con Jer 23:5--6 y 31:35--37, presentes en ambas formas del texto.
\end{notalam}
\end{column}
\begin{column}{0.46\textwidth}\centering
\begin{tikzpicture}[node distance=5mm]
\node[bloqueazul,text width=4.2cm] (r) {Renuevo justo de David\\{\normalfont\tiny Jer 23:5}};
\node[bloqueoro,text width=1.85cm,below left=7mm and -1.9cm of r] (j) {\textbf{Judá}\\«será salvo»};
\node[bloque,text width=1.85cm,below right=7mm and -1.9cm of r] (i) {\textbf{Israel}\\«habitará confiado»};
\draw[flecha] (r) -- (j); \draw[flecha] (r) -- (i);
\end{tikzpicture}
\elabpropia
\vspace{2pt}
\footnotesize
\begin{itemize}
\item \emph{YHWH tsidqenu} frente a Sedequías (\emph{Tsidqiyahu}).
\item Jer 33:20--21: pacto con David como el pacto con el día y la noche.
\item Jer 33:24--26: «las dos familias» no serán desechadas.
\end{itemize}
\end{column}
\end{columns}
\end{frame}

\begin{frame}{Ezequiel: dos palos}
\fotolado{panel_ez37.jpg}{Elaboración propia: Ez 37:22 en texto hebreo consonántico (BHS).}
\begin{columns}[T,onlytextwidth]
\begin{column}{0.57\textwidth}
\begin{cita}Los haré una nación en la tierra, en los montes de Israel, y un rey será a todos ellos por rey; y nunca más serán dos naciones, ni nunca más serán divididos en dos reinos. [\ldots] Mi siervo David será rey sobre ellos.\fuente{(Ez 37:22, 24)}\end{cita}
\centering
\begin{tikzpicture}[node distance=3mm]
\node[bloqueoro,text width=2.2cm] (j) {palo «para Judá»};
\node[bloque,text width=2.2cm,below=of j] (e) {palo «para José, palo de Efraín»};
\node[bloqueazul,text width=2.5cm,anchor=west] (u) at ($(j.east)!0.5!(e.east)+(1.0,0)$) {una nación, un rey: «David»\\{\normalfont\tiny pacto perpetuo, 37:26--28}};
\draw[flecha] (j.east) -- (u.west); \draw[flecha] (e.east) -- (u.west);
\end{tikzpicture}
\notafig{Elaboración propia a partir de Ez 37:16--28. El oráculo deshace 1~R 12.}
\end{column}
\begin{column}{0.41\textwidth}\end{column}
\end{columns}
\end{frame}

\begin{frame}{Amós y Oseas}
\begin{columns}[T,onlytextwidth]
\begin{column}{0.48\textwidth}
\textbf{\textcolor{cafe}{Amós, en Betel}}
\begin{cita}[fontupper=\scriptsize\itshape]En aquel día yo levantaré la cabaña caída de David [\ldots], para que posean el resto de Edom y todas las naciones sobre las cuales es invocado mi nombre.\fuente{(Am 9:11--12)}\end{cita}
\footnotesize Jacobo lo cita en Hch 15:16--17 según la LXX: «para que el resto de los hombres busque al Señor».
\begin{notalam}[TM y LXX]
TM \emph{yirshú} «posean» / \emph{Edom}; LXX «busquen» / «hombres» (\emph{adam}): consonantes casi iguales.
\end{notalam}
\end{column}
\begin{column}{0.48\textwidth}
\textbf{\textcolor{cafe}{Oseas, profeta del norte}}
\begin{cita}[fontupper=\scriptsize\itshape]Después volverán los hijos de Israel, y buscarán a YHWH su Dios, y a David su rey.\fuente{(Os 3:5)}\end{cita}
\footnotesize El norte regresará a la casa que rechazó en Siquem. «Se congregarán los hijos de Judá y de Israel, y nombrarán un solo jefe» (Os 1:11 [2:2]).\par\smallskip
\begin{clave}ONE FOR ISRAEL (2025) cita Os 3:4--5 entre los textos que anuncian un regreso futuro del pueblo judío al Mesías.\end{clave}
\end{column}
\end{columns}
\end{frame}

\begin{frame}{Lucas y Pablo: el hijo de David}
\fotolado{panel_lc1.jpg}{Elaboración propia: Lc 1:32--33 en escritura uncial (texto NA28).}
\begin{columns}[T,onlytextwidth]
\begin{column}{0.57\textwidth}
\begin{cita}El Señor Dios le dará el trono de David su padre; y reinará sobre la casa de Jacob para siempre.\fuente{(Lc 1:32--33)}\end{cita}
\scriptsize
\begin{itemize}
\item \textbf{Ro 1:3--4}: \grc{ἐκ σπέρματος Δαυὶδ κατὰ σάρκα}, declarado Hijo de Dios por la resurrección (2~S 7:12--14).
\item \textbf{Ro 15:12}: la raíz de Isaí; las naciones esperarán en él.
\item \textbf{Ro 11:1--4}: «¿Ha desechado Dios a su pueblo?» Pablo, benjaminita; Elías, ejemplo del norte.
\item \textbf{Ro 11:26}: «Vendrá de \textbf{Sion} el Libertador».
\end{itemize}
\precision{2~Ti 2:8 dice lo mismo; su autoría se discute, pero Ro 1:3--4 basta.}
\end{column}
\begin{column}{0.41\textwidth}\end{column}
\end{columns}
\end{frame}

% ===================================================
\section{Síntesis}
\begin{frame}{Síntesis}
\begin{tcolorbox}[colback=piedra,colframe=ladrillo,boxrule=1pt,arc=3pt,fontupper=\footnotesize]
\textbf{\textcolor{cafe}{Israel de David a la división:}} un pueblo unido por pacto bajo un rey a quien Dios prometió una casa «para siempre» y un templo abierto a las naciones; después, un reino partido por la idolatría y la opresión, sometido a un juicio que tiene término: «no para siempre» (1~R 11:39). Los profetas esperan un hijo de David bajo el cual Efraín y Judá vuelvan a ser una nación (Ez 37:22).
\end{tcolorbox}
\vspace{2pt}
\centering\footnotesize\renewcommand{\arraystretch}{1.25}
\begin{tabular}{>{\RaggedRight\arraybackslash}p{2.6cm}>{\RaggedRight\arraybackslash}p{8.6cm}}
\toprule
\color{ladrillo}\bfseries Eje & \color{ladrillo}\bfseries Lo que deja esta etapa\\\midrule
Unidad frágil & Pacto de Hebrón; tensión norte--sur (2~S 5:3; 19:41--20:1)\\
Vocación & Arca en Sion; templo para los pueblos; Sal 72:17\\
Restauración & 2~S 7:14--16; 1~R 8:46--53; 11:36--39; Sal 89\\
Proyección & Is 11; Jer 23; Ez 37; Am 9; Os 3:5; Ro 1:3--4; 11:1, 26\\
\bottomrule
\end{tabular}
\elabpropia
\end{frame}

\begin{frame}{¿Cómo volverán a ser uno?}
\begin{cita}Júntalos luego el uno con el otro, para que sean uno solo, y serán uno solo en tu mano.\fuente{(Ez 37:17)}\end{cita}
\centering
\begin{tikzpicture}
\node[bloqueazul,text width=1.9cm,minimum height=1.2cm] (d) at (0,0) {División\\{\normalfont\tiny 1~R 12, c. 931/930}};
\node[bloque,text width=2.5cm] (n1) at (3.4,0.85) {\textbf{Israel (norte)}\\{\tiny capital Samaria}};
\node[bloque,text width=2.5cm] (n2) at (6.9,0.85) {Cae ante \textbf{Asiria}\\{\tiny 722 a.\,e.\,c.}};
\node[bloqueoro,text width=2.5cm] (s1) at (3.4,-0.85) {\textbf{Judá (sur)}\\{\tiny casa de David}};
\node[bloqueoro,text width=2.5cm] (s2) at (6.9,-0.85) {Cae ante \textbf{Babilonia}\\{\tiny 586 a.\,e.\,c.} y vuelve};
\node[bloqueazul,text width=1.9cm,minimum height=1.2cm] (u) at (10.4,0) {«Uno solo en tu mano»\\{\normalfont\tiny Ez 37:17}};
\draw[flecha] (d) -- (n1); \draw[flecha] (d) -- (s1);
\draw[flecha] (n1) -- (n2); \draw[flecha] (s1) -- (s2);
\draw[flecha,dashed] (n2) -- (u); \draw[flecha,dashed] (s2) -- (u);
\end{tikzpicture}
\elabpropia
\vspace{4pt}
\begin{tcolorbox}[colback=cafe,colframe=cafeosc,coltext=piedra,boxrule=0.8pt,arc=3pt,fontupper=\footnotesize]
\textbf{Estudio 5.} Los reyes y los cautiverios: reino dividido, Asiria y Babilonia, profetas y promesas de nuevo pacto y reunificación (Jer 31; Ez 36--37).
\end{tcolorbox}
\end{frame}

\begin{frame}{Para la reflexión}
\small
\begin{enumerate}\setlength{\itemsep}{6pt}
\item David quiso construir una casa para Dios, y Dios le prometió una casa a él (2~S 7:5, 11). ¿Qué dice ese giro sobre cómo Dios cumple su propósito?
\item ¿Cómo se sostiene una promesa «para siempre» cuando los hechos parecen negarla (Sal 89)?
\item Salomón oró por el extranjero y levantó altares a dioses extranjeros. ¿Qué enseña esa contradicción?
\item ¿Qué parte de la división fue responsabilidad de Judá y qué parte de Israel (norte)?
\item Pablo, benjaminita, toma su ejemplo de remanente del reino del norte (Ro 11:1--4). ¿Qué sugiere sobre «todo Israel»?
\end{enumerate}
\end{frame}

\begin{frame}{Glosario}
\begin{columns}[T,onlytextwidth]
\begin{column}{0.48\textwidth}\small
\textbf{Ad olam} (\heb{עַד־עוֹלָם}): «para siempre» (2~S 7:16).\par\smallskip
\textbf{Báyit} (\heb{בַּיִת}): casa: palacio, templo, dinastía.\par\smallskip
\textbf{Bet David} (\heb{בֵּית דָּוִד}): casa de David; \emph{bytdwd} en Tel Dan.\par\smallskip
\textbf{Efraín} (\heb{אֶפְרַיִם}): tribu central; nombre del reino del norte.\par\smallskip
\textbf{Jésed} (\heb{חֶסֶד}): misericordia leal al pacto.\par\smallskip
\textbf{Kun} (\heb{כּוּן}): establecer; \emph{najón}, \emph{Yajín}.
\end{column}
\begin{column}{0.48\textwidth}\small
\textbf{Mas} (\heb{מַס}): leva de trabajo forzado (Éx 1:11; 1~R 5:13).\par\smallskip
\textbf{Mashíaj} (\heb{מָשִׁיחַ}): ungido; \grc{Χριστός}.\par\smallskip
\textbf{Nir} (\heb{נִיר}): lámpara de David (1~R 11:36).\par\smallskip
\textbf{Shuv} (\heb{שׁוּב}): volver, convertirse (1~R 8:47).\par\smallskip
\textbf{Sukkat David} (\heb{סֻכַּת דָּוִיד}): cabaña de David (Am 9:11).\par\smallskip
\textbf{Tsémaj} (\heb{צֶמַח}): renuevo (Jer 23:5; Amidá 15).
\end{column}
\end{columns}
\end{frame}

\begin{frame}{Referencias}
\begin{columns}[t,onlytextwidth]
\begin{column}{0.49\textwidth}\fontsize{4.5}{5.2}\selectfont
\refapa{Alter, R. (1999). \emph{The David story}. W.~W. Norton.}
\refapa{Ben-Yosef, E. (2019). The architectural bias in current biblical archaeology. \emph{Vetus Testamentum, 69}(3), 361--387.}
\refapa{\emph{Biblia Hebraica Stuttgartensia} (5.ª ed.). (1997). Deutsche Bibelgesellschaft.}
\refapa{Biran, A., \& Naveh, J. (1993). An Aramaic stele fragment from Tel Dan. \emph{Israel Exploration Journal, 43}(2--3), 81--98.}
\refapa{Biran, A., \& Naveh, J. (1995). The Tel Dan inscription: A new fragment. \emph{Israel Exploration Journal, 45}(1), 1--18.}
\refapa{Broshi, M., \& Finkelstein, I. (1992). The population of Palestine in Iron Age II. \emph{BASOR, 287}, 47--60.}
\refapa{Charlesworth, J.~H. (Ed.). (1983--1985). \emph{The Old Testament pseudepigrapha}. Doubleday.}
\refapa{Chilton, B.~D. (1987). \emph{The Isaiah Targum}. Michael Glazier.}
\refapa{Cogan, M. (2001). \emph{1 Kings} (Anchor Bible 10). Doubleday.}
\refapa{Elbogen, I. (1993). \emph{Jewish liturgy: A comprehensive history}. Jewish Publication Society.}
\refapa{Epstein, I. (Ed.). (1935--1948). \emph{The Babylonian Talmud}. Soncino Press.}
\refapa{Finkelstein, I. (1996). The archaeology of the United Monarchy: An alternative view. \emph{Levant, 28}, 177--187.}
\refapa{Finkelstein, I., \& Fantalkin, A. (2012). Khirbet Qeiyafa: An unsensational archaeological and historical interpretation. \emph{Tel Aviv, 39}(1), 38--63.}
\refapa{Finkelstein, I., Herzog, Z., Singer-Avitz, L., \& Ussishkin, D. (2007). Has King David's palace in Jerusalem been found? \emph{Tel Aviv, 34}(2), 142--164.}
\refapa{Finkelstein, I., Na'aman, N., \& Römer, T. (2019). Restoring line 31 in the Mesha Stele. \emph{Tel Aviv, 46}(1), 3--11.}
\refapa{Finkelstein, I., \& Silberman, N.~A. (2001). \emph{The Bible unearthed}. Free Press.}
\refapa{Finkelstein, I., \& Silberman, N.~A. (2006). \emph{David and Solomon}. Free Press.}
\refapa{Garfinkel, Y., Kreimerman, I., \& Zilberg, P. (2016). \emph{Debating Khirbet Qeiyafa}. Israel Exploration Society.}
\end{column}
\begin{column}{0.49\textwidth}\fontsize{4.5}{5.2}\selectfont
\refapa{Hays, R.~B. (1989). \emph{Echoes of Scripture in the letters of Paul}. Yale University Press.}
\refapa{Josefo, F. (1987). \emph{The works of Josephus} (W. Whiston, Trad.). Hendrickson.}
\refapa{Kitchen, K.~A. (1986). \emph{The Third Intermediate Period in Egypt} (2.ª ed.). Aris \& Phillips.}
\refapa{Kitchen, K.~A. (2003). \emph{On the reliability of the Old Testament}. Eerdmans.}
\refapa{Lemaire, A. (1994). «House of David» restored in Moabite inscription. \emph{Biblical Archaeology Review, 20}(3), 30--37.}
\refapa{Mazar, A. (2010). Archaeology and the biblical narrative: The case of the United Monarchy. En R.~G. Kratz \& H. Spieckermann (Eds.), \emph{One God -- one cult -- one nation}. De Gruyter.}
\refapa{Mazar, E. (2006). Did I find King David's palace? \emph{Biblical Archaeology Review, 32}(1).}
\refapa{McCarter, P.~K. (1984). \emph{II Samuel} (Anchor Bible 9). Doubleday.}
\refapa{Mendenhall, G.~E. (1958). The census lists of Numbers 1 and 26. \emph{Journal of Biblical Literature, 77}(1), 52--66.}
\refapa{Monson, J. (2000). The new \textquoteleft Ain Dara temple: Closest Solomonic parallel. \emph{Biblical Archaeology Review, 26}(3).}
\refapa{\emph{Novum Testamentum Graece} (28.ª ed.). (2012). Deutsche Bibelgesellschaft.}
\refapa{ONE FOR ISRAEL Ministry. (2025, 3 de julio). \emph{What does «all Israel will be saved» mean in Romans 11?} \url{https://www.oneforisrael.org/what-does-all-israel-will-be-saved-mean-in-romans-11/}}
\refapa{Rahlfs, A., \& Hanhart, R. (Eds.). (2006). \emph{Septuaginta} (ed. rev.). Deutsche Bibelgesellschaft.}
\refapa{Reich, R., \& Shukron, E. (1999). Light at the end of the tunnel. \emph{Biblical Archaeology Review, 25}(1).}
\refapa{Sacks, J. (Trad.). (2009). \emph{The Koren Siddur}. Koren Publishers.}
\refapa{Thiele, E.~R. (1983). \emph{The mysterious numbers of the Hebrew kings} (ed. rev.). Zondervan.}
\refapa{Vermes, G. (2004). \emph{The complete Dead Sea Scrolls in English} (ed. rev.). Penguin.}
\refapa{Wright, N.~T. (2013). \emph{Paul and the faithfulness of God}. Fortress Press.}
\refapa{Yadin, Y. (1958). Solomon's city wall and gate at Gezer. \emph{Israel Exploration Journal, 8}(2), 80--86.}
\end{column}
\end{columns}
\end{frame}
\end{document}
